\documentclass[tikz,border=6pt]{standalone}
\usepackage{ctex}
\usepackage{amsmath}
\usepackage{pgfplots}
\pgfplotsset{compat=1.18}
\usetikzlibrary{calc,angles,quotes,arrows.meta,positioning,decorations.markings}
\begin{document}
\begin{tikzpicture}[scale=3.4,>=Stealth]
  \def\alp{68}
  \def\bet{23}

  % 坐标轴
  \draw[->,gray!70] (-1.25,0) -- (1.35,0) node[below right,gray!80] {$x$};
  \draw[->,gray!70] (0,-0.45) -- (0,1.35) node[above left,gray!80] {$y$};

  % 单位圆
  \draw[thick,blue!55] (0,0) circle (1);

  \coordinate (O) at (0,0);
  \coordinate (P) at (\alp:1);
  \coordinate (Q) at (\bet:1);

  % 半径
  \draw[thick,red!70] (O) -- (P);
  \draw[thick,teal!75] (O) -- (Q);

  % 连接 P,Q 的弦（两点距离法的弦）
  \draw[very thick,violet!75,dashed] (P) -- (Q);

  % 角度弧：alpha, beta, 以及夹角 alpha-beta
  \draw[->,thick,red!70] (0.30,0) arc (0:\alp:0.30);
  \node[red!70] at (\alp*0.5:0.20) {$\alpha$};

  \draw[->,thick,teal!75] (0.62,0) arc (0:\bet:0.62);
  \node[teal!75] at (\bet*0.5:0.72) {$\beta$};

  % 夹角 alpha - beta（在两半径之间，靠外侧标注）
  \draw[->,thick,violet!75] (\bet:0.45) arc (\bet:\alp:0.45);
  \node[violet!75] at ({(\alp+\bet)*0.5}:0.33) {$\alpha-\beta$};

  % 点标注
  \fill[red!70] (P) circle (0.018);
  \node[above right] at (P) {$P_1(\cos\alpha,\,\sin\alpha)$};
  \fill[teal!75] (Q) circle (0.018);
  \node[right] at (Q) {$P_2(\cos\beta,\,\sin\beta)$};

  % 半透明注释框（放右下空白区，避开圆与弦）
  \node[fill=white,fill opacity=0.80,text opacity=1,rounded corners,
        draw=violet!45,inner sep=4pt,align=center,font=\small]
        at (0.74,-0.26)
        {两点距离法：\\[1pt]$|P_1P_2|^2=2-2\cos(\alpha-\beta)$\\[1pt]
         $\Rightarrow\cos(\alpha-\beta)=$\\[1pt]
         $\cos\alpha\cos\beta+\sin\alpha\sin\beta$};
\end{tikzpicture}
\end{document}
